% TOP_PROBLEM: {"id": 3700025, "title": "Two one-points in the unit disk", "category_id": 8, "category": {"id": 8, "name": "analysis", "display_name": "Analysis", "description": "Limits, continuity, calculus, and function theory.", "slug": "analysis", "order_index": 8, "created_at": "2026-07-31T15:26:25.671Z"}, "status": "partially_solved", "problem_number": "AMR-036-0025", "set_id": 13}
% FIRST_POSTED: 2026-10-01
% ENTRY_KIND: statement_only
% UnsolvedMath ID 3700025; source catalogue code AMR-036-0025
% !TeX program = lualatex
% Definition and sources only; this file is not an LLM solution attempt.
\documentclass[11pt,a4paper]{article}
\usepackage[margin=25mm]{geometry}
\usepackage{fontspec}
\setmainfont{lmroman10-regular.otf}[BoldFont=lmroman10-bold.otf,
 ItalicFont=lmroman10-italic.otf,BoldItalicFont=lmroman10-bolditalic.otf]
\usepackage{amsmath,amssymb,mathtools}
\usepackage{unicode-math}
\setmathfont{latinmodern-math.otf}
\AtBeginDocument{\let\setminus\smallsetminus}
\usepackage{xurl,hyperref}
\hypersetup{colorlinks=true,urlcolor=blue}
\setcounter{secnumdepth}{0}
\setlength{\parindent}{0pt}
\setlength{\parskip}{5pt}
\setlength{\emergencystretch}{3em}
\begin{document}
\section*{Two one-points in the unit disk}\label{3700025}
{\small UnsolvedMath ID 3700025 · AMR-036-0025 · Analysis · AMR Open Problem Lists}
\subsection{Definitions and mathematical statement}
Let $\mathcal F$ be the class of holomorphic functions
$f$ on the unit disk $\mathbb D$ satisfying
$f(0)=0$, $f'(0)\ne0$, and $f(z)\ne0$ for every
$z\in\mathbb D\setminus\{0\}$. Require that $f-1$
have exactly two zeros in $\mathbb D$, counted with
multiplicity, and denote them by $z_1,z_2$.
A single double zero of $f-1$ is allowed and is
represented by $z_1=z_2$.

Define the two extremal quantities
\[
 R_*=\inf_{f\in\mathcal F}\max\{|z_1|,|z_2|\},
 \qquad
 D_*=\sup_{f\in\mathcal F}|f'(0)|.
\]
Determine their sharp values, decide whether the infimum
and supremum are attained, and characterize the
corresponding extremal functions when they exist.
If attainment fails, describe admissible functions approaching
the sharp value.

The zero at the origin is fixed and simple, so arbitrary
precomposition by a disk automorphism is not a permitted
normalization. Rotations of the variable remain allowed
and preserve both quantities. The radius objective is
Euclidean and centered at the origin, unlike a freely
centered hyperbolic enclosing radius.
The derivative objective is a separate optimization;
the formulation does not assume that one function must
extremize both. There are no other restrictions on the
range of $f$, its behavior at the boundary, or the
multiplicity of its critical points.
\subsection{Short English statement}
For a disk function with one simple zero at the origin and two one-points, find the smallest enclosing centered radius and the largest derivative.
\subsection{Sources}
\begin{itemize}
\item[S1] ULAM.AI and the UnsolvedMath Contributors, supplied problems.json, record 3700025 (AMR-036-0025), AMR Open Problem Lists. Catalogue identity and source transcription; CC BY 4.0.
\url{https://huggingface.co/datasets/ulamai/UnsolvedMath}.
\item[S2] Eremenko - My favorite unsolved problems, Goldberg relatives, Problem 1. Original source indexed by AMR-036-0025.
\url{https://www.math.purdue.edu/~eremenko/uns1.html}.
\item[S3] A. Eremenko, Goldberg's constant and its relatives, Problem 1. Author-maintained primary problem note.
\url{https://www.math.purdue.edu/~eremenko/dvi/goldconst.pdf}.
\end{itemize}
\end{document}
